\clearpage

\section{完整程序代码}

\label{app:complete_code}

本附录逐文件收录40份程序与依赖文件，共8360行，覆盖数据读取、预测、优化、候选筛选、对照实验、图表、工作簿及检查。代码由源文件全文读入，保留注释与原始行号；长行折行时沿用同一行号。

\subsection*{运行环境与复现顺序}

Python 数值程序依赖 NumPy、SciPy、Matplotlib 和 OpenPyXL，LP 由 SciPy 的 HiGHS 接口求解。\path{code/requirements.txt}保留原交付环境清单；最终候选及消融运行记录使用 Python 3.11.15、NumPy 2.4.6、SciPy 1.17.1。Node.js 工作簿导出另依赖\path{@oai/artifact-tool} 2.8.59；运行执行目录副本时，设置\path{CUMCM_PROJECT_ROOT}为项目根目录，并确保依赖可被 Node.js 解析。PowerShell 入口用于原运行环境；论文使用 XeLaTeX 编译，代码采用 Menlo 与宋体字体。

复现按“附件与时间映射—一月配置冻结—母模型与外层筛选—结果发布及核对”推进。\path{run_all.ps1}及\path{run_revision.ps1}分别对应原版和母模型；最终滚动方案另由候选与时点实验入口生成，之后执行\path{publish_final_strategy.py}、工作簿导出与核对。\path{make_compact_paper_tables.py}恢复紧凑表排布。冻结配置、实测附件、空白模板及执行数组按\ref{tab:files}所述目录组织。

原版诊断按保存路径动态加载冻结实现，因此一并列出该文件；工作簿执行目录副本也完整列出。旧打包及诊断脚本保留当时版本的说明，其状态文字不作为当前最终方案的结果。各源码按下列相对路径保存，保持目录结构；第三方库按依赖安装，历史备份中的重复文件不再逐份列入。

\subsection*{文件与页码索引}

\begingroup\small\setlength{\tabcolsep}{4pt}

\begin{longtable}{r p{0.44\linewidth} p{0.30\linewidth} r r}

\toprule 编号 & 文件 & 用途 & 行数 & 页码\\\midrule\endfirsthead

\toprule 编号 & 文件 & 用途 & 行数 & 页码\\\midrule\endhead

\midrule\multicolumn{5}{r}{续下页}\\\endfoot

\bottomrule\endlastfoot

1 & \path{code/requirements.txt} & 原交付环境依赖清单 & 5 & \pageref{code:source:01}\\

2 & \path{code/data_forecast.py} & 附件读取与因果预测 & 344 & \pageref{code:source:02}\\

3 & \path{code/time_index.py} & 物理时段与附件索引映射 & 196 & \pageref{code:source:03}\\

4 & \path{code/optimize_dispatch.py} & 确定性模型与原版场景调度 & 387 & \pageref{code:source:04}\\

5 & \path{code/dispatch_v2.py} & 统一母模型、控制器与结算 & 309 & \pageref{code:source:05}\\

6 & \path{code/calibrate_v2.py} & 一月配置校准与冻结 & 78 & \pageref{code:source:06}\\

7 & \path{code/run_main_v2.py} & 四组母模型计算入口 & 44 & \pageref{code:source:07}\\

8 & \path{code/outer_candidate_selector.py} & 候选生成与执行回放评分 & 246 & \pageref{code:source:08}\\

9 & \path{code/run_outer_candidate_experiment.py} & 全时点筛选实验 & 286 & \pageref{code:source:09}\\

10 & \path{code/outer_timing_policy.py} & 筛选启用时点策略 & 175 & \pageref{code:source:10}\\

11 & \path{code/run_outer_timing_ablation.py} & 筛选时点消融与最终选型 & 261 & \pageref{code:source:11}\\

12 & \path{code/run_experiments.py} & 原版模型对照实验 & 160 & \pageref{code:source:12}\\

13 & \path{code/experiments_v2.py} & 母模型风险、备用与反馈对照 & 373 & \pageref{code:source:13}\\

14 & \path{code/diagnose_execution_gap.py} & 冻结原版规划与执行诊断 & 199 & \pageref{code:source:14}\\

15 & \path{code/diagnose_inner_v2.py} & 母模型执行器及续接诊断 & 226 & \pageref{code:source:15}\\

16 & \path{code/scenario_convergence.py} & 历史场景压缩与收敛记录 & 86 & \pageref{code:source:16}\\

17 & \path{code/make_paper_assets.py} & 原版结果图表生成 & 231 & \pageref{code:source:17}\\

18 & \path{code/make_paper_assets_v2.py} & 母模型结果图表生成 & 213 & \pageref{code:source:18}\\

19 & \path{code/make_question_tables.py} & 分问指定结果表生成 & 175 & \pageref{code:source:19}\\

20 & \path{code/make_visual_atlas.py} & 调度与诊断图生成 & 464 & \pageref{code:source:20}\\

21 & \path{code/publish_final_strategy.py} & 最终方案指标与图表发布 & 248 & \pageref{code:source:21}\\

22 & \path{code/make_compact_paper_tables.py} & 紧凑指定日期表生成 & 107 & \pageref{code:source:22}\\

23 & \path{code/export_workbooks.mjs} & 原版与母模型工作簿填写 & 261 & \pageref{code:source:23}\\

24 & \path{code/export_final_workbooks.mjs} & 最终滚动工作簿填写 & 261 & \pageref{code:source:24}\\

25 & \path{code/verify_workbooks.py} & 原版工作簿检查 & 187 & \pageref{code:source:25}\\

26 & \path{code/verify_final_workbooks.py} & 最终工作簿检查 & 187 & \pageref{code:source:26}\\

27 & \path{code/check_dispatch_v2.py} & 母模型数值与信息检查 & 57 & \pageref{code:source:27}\\

28 & \path{code/check_outer_candidate_selector.py} & 候选回放单项检查 & 312 & \pageref{code:source:28}\\

29 & \path{code/check_outer_timing_policy.py} & 筛选时点策略检查 & 261 & \pageref{code:source:29}\\

30 & \path{code/check_outer_timing_ablation.py} & 筛选时点消融检查 & 359 & \pageref{code:source:30}\\

31 & \path{code/check_final_strategy_delivery.py} & 最终执行与交付检查 & 449 & \pageref{code:source:31}\\

32 & \path{code/check_manuscript_sources.py} & 原分问稿来源与结构检查 & 61 & \pageref{code:source:32}\\

33 & \path{code/run_all.ps1} & 原版完整运行入口 & 46 & \pageref{code:source:33}\\

34 & \path{code/run_revision.ps1} & 母模型修订复算入口 & 153 & \pageref{code:source:34}\\

35 & \path{code/package_revision.py} & 原分问稿打包工具 & 65 & \pageref{code:source:35}\\

36 & \path{versions/v1_before_settlement_revision/code/optimize_dispatch.py} & 原版诊断动态加载的冻结实现 & 387 & \pageref{code:source:36}\\

37 & \path{compile.sh} & XeLaTeX 编译入口 & 36 & \pageref{code:source:37}\\

38 & \path{outputs/final_strategy/runtime/export_final_workbooks.mjs} & 工作簿导出的执行目录副本 & 261 & \pageref{code:source:38}\\

39 & \path{outputs/logic_formula_review/verify_revision.py} & 公式与版式专项审计 & 88 & \pageref{code:source:39}\\

40 & \path{outputs/code_appendix/build_code_appendix.py} & 本代码附录生成器 & 116 & \pageref{code:source:40}\\

\end{longtable}\endgroup

\clearpage

\begingroup

\newfontfamily\appendixcodefont{Menlo}[NFSSFamily=appendixcode]

\setCJKfamilyfont{appendixcjk}{Songti SC}

\fvset{fontfamily=appendixcode,fontsize=\fontsize{8.5}{10.5}\selectfont,baselinestretch=1,formatcom=\CJKfamily{appendixcjk},numbers=left,numbersep=6pt,xleftmargin=2.8em,frame=lines,framesep=5pt,rulecolor=\color{black!25},breaklines=true,breakanywhere=true,breakautoindent=true,breakindent=1em,breaksymbolleft=\tiny\ensuremath{\hookrightarrow},tabsize=4,showspaces=false,showtabs=false}

\renewcommand{\theFancyVerbLine}{\scriptsize\color{black!50}\arabic{FancyVerbLine}}

\Needspace{7\baselineskip}

\subsection{原交付环境依赖清单}

\label{code:source:01}

{\small\noindent\path{code/requirements.txt}\quad（5行）\par}

\VerbatimInput{code/requirements.txt}

\label{code:source:01:end}

\Needspace{7\baselineskip}

\subsection{附件读取与因果预测}

\label{code:source:02}

{\small\noindent\path{code/data_forecast.py}\quad（344行）\par}

\VerbatimInput{code/data_forecast.py}

\label{code:source:02:end}

\Needspace{7\baselineskip}

\subsection{物理时段与附件索引映射}

\label{code:source:03}

{\small\noindent\path{code/time_index.py}\quad（196行）\par}

\VerbatimInput{code/time_index.py}

\label{code:source:03:end}

\Needspace{7\baselineskip}

\subsection{确定性模型与原版场景调度}

\label{code:source:04}

{\small\noindent\path{code/optimize_dispatch.py}\quad（387行）\par}

\VerbatimInput{code/optimize_dispatch.py}

\label{code:source:04:end}

\Needspace{7\baselineskip}

\subsection{统一母模型、控制器与结算}

\label{code:source:05}

{\small\noindent\path{code/dispatch_v2.py}\quad（309行）\par}

\VerbatimInput{code/dispatch_v2.py}

\label{code:source:05:end}

\Needspace{7\baselineskip}

\subsection{一月配置校准与冻结}

\label{code:source:06}

{\small\noindent\path{code/calibrate_v2.py}\quad（78行）\par}

\VerbatimInput{code/calibrate_v2.py}

\label{code:source:06:end}

\Needspace{7\baselineskip}

\subsection{四组母模型计算入口}

\label{code:source:07}

{\small\noindent\path{code/run_main_v2.py}\quad（44行）\par}

\VerbatimInput{code/run_main_v2.py}

\label{code:source:07:end}

\Needspace{7\baselineskip}

\subsection{候选生成与执行回放评分}

\label{code:source:08}

{\small\noindent\path{code/outer_candidate_selector.py}\quad（246行）\par}

\VerbatimInput{code/outer_candidate_selector.py}

\label{code:source:08:end}

\Needspace{7\baselineskip}

\subsection{全时点筛选实验}

\label{code:source:09}

{\small\noindent\path{code/run_outer_candidate_experiment.py}\quad（286行）\par}

\VerbatimInput{code/run_outer_candidate_experiment.py}

\label{code:source:09:end}

\Needspace{7\baselineskip}

\subsection{筛选启用时点策略}

\label{code:source:10}

{\small\noindent\path{code/outer_timing_policy.py}\quad（175行）\par}

\VerbatimInput{code/outer_timing_policy.py}

\label{code:source:10:end}

\Needspace{7\baselineskip}

\subsection{筛选时点消融与最终选型}

\label{code:source:11}

{\small\noindent\path{code/run_outer_timing_ablation.py}\quad（261行）\par}

\VerbatimInput{code/run_outer_timing_ablation.py}

\label{code:source:11:end}

\Needspace{7\baselineskip}

\subsection{原版模型对照实验}

\label{code:source:12}

{\small\noindent\path{code/run_experiments.py}\quad（160行）\par}

\VerbatimInput{code/run_experiments.py}

\label{code:source:12:end}

\Needspace{7\baselineskip}

\subsection{母模型风险、备用与反馈对照}

\label{code:source:13}

{\small\noindent\path{code/experiments_v2.py}\quad（373行）\par}

\VerbatimInput{code/experiments_v2.py}

\label{code:source:13:end}

\Needspace{7\baselineskip}

\subsection{冻结原版规划与执行诊断}

\label{code:source:14}

{\small\noindent\path{code/diagnose_execution_gap.py}\quad（199行）\par}

\VerbatimInput{code/diagnose_execution_gap.py}

\label{code:source:14:end}

\Needspace{7\baselineskip}

\subsection{母模型执行器及续接诊断}

\label{code:source:15}

{\small\noindent\path{code/diagnose_inner_v2.py}\quad（226行）\par}

\VerbatimInput{code/diagnose_inner_v2.py}

\label{code:source:15:end}

\Needspace{7\baselineskip}

\subsection{历史场景压缩与收敛记录}

\label{code:source:16}

{\small\noindent\path{code/scenario_convergence.py}\quad（86行）\par}

\VerbatimInput{code/scenario_convergence.py}

\label{code:source:16:end}

\Needspace{7\baselineskip}

\subsection{原版结果图表生成}

\label{code:source:17}

{\small\noindent\path{code/make_paper_assets.py}\quad（231行）\par}

\VerbatimInput{code/make_paper_assets.py}

\label{code:source:17:end}

\Needspace{7\baselineskip}

\subsection{母模型结果图表生成}

\label{code:source:18}

{\small\noindent\path{code/make_paper_assets_v2.py}\quad（213行）\par}

\VerbatimInput{code/make_paper_assets_v2.py}

\label{code:source:18:end}

\Needspace{7\baselineskip}

\subsection{分问指定结果表生成}

\label{code:source:19}

{\small\noindent\path{code/make_question_tables.py}\quad（175行）\par}

\VerbatimInput{code/make_question_tables.py}

\label{code:source:19:end}

\Needspace{7\baselineskip}

\subsection{调度与诊断图生成}

\label{code:source:20}

{\small\noindent\path{code/make_visual_atlas.py}\quad（464行）\par}

\VerbatimInput{code/make_visual_atlas.py}

\label{code:source:20:end}

\Needspace{7\baselineskip}

\subsection{最终方案指标与图表发布}

\label{code:source:21}

{\small\noindent\path{code/publish_final_strategy.py}\quad（248行）\par}

\VerbatimInput{code/publish_final_strategy.py}

\label{code:source:21:end}

\Needspace{7\baselineskip}

\subsection{紧凑指定日期表生成}

\label{code:source:22}

{\small\noindent\path{code/make_compact_paper_tables.py}\quad（107行）\par}

\VerbatimInput{code/make_compact_paper_tables.py}

\label{code:source:22:end}

\Needspace{7\baselineskip}

\subsection{原版与母模型工作簿填写}

\label{code:source:23}

{\small\noindent\path{code/export_workbooks.mjs}\quad（261行）\par}

\VerbatimInput{code/export_workbooks.mjs}

\label{code:source:23:end}

\Needspace{7\baselineskip}

\subsection{最终滚动工作簿填写}

\label{code:source:24}

{\small\noindent\path{code/export_final_workbooks.mjs}\quad（261行）\par}

\VerbatimInput{code/export_final_workbooks.mjs}

\label{code:source:24:end}

\Needspace{7\baselineskip}

\subsection{原版工作簿检查}

\label{code:source:25}

{\small\noindent\path{code/verify_workbooks.py}\quad（187行）\par}

\VerbatimInput{code/verify_workbooks.py}

\label{code:source:25:end}

\Needspace{7\baselineskip}

\subsection{最终工作簿检查}

\label{code:source:26}

{\small\noindent\path{code/verify_final_workbooks.py}\quad（187行）\par}

\VerbatimInput{code/verify_final_workbooks.py}

\label{code:source:26:end}

\Needspace{7\baselineskip}

\subsection{母模型数值与信息检查}

\label{code:source:27}

{\small\noindent\path{code/check_dispatch_v2.py}\quad（57行）\par}

\VerbatimInput{code/check_dispatch_v2.py}

\label{code:source:27:end}

\Needspace{7\baselineskip}

\subsection{候选回放单项检查}

\label{code:source:28}

{\small\noindent\path{code/check_outer_candidate_selector.py}\quad（312行）\par}

\VerbatimInput{code/check_outer_candidate_selector.py}

\label{code:source:28:end}

\Needspace{7\baselineskip}

\subsection{筛选时点策略检查}

\label{code:source:29}

{\small\noindent\path{code/check_outer_timing_policy.py}\quad（261行）\par}

\VerbatimInput{code/check_outer_timing_policy.py}

\label{code:source:29:end}

\Needspace{7\baselineskip}

\subsection{筛选时点消融检查}

\label{code:source:30}

{\small\noindent\path{code/check_outer_timing_ablation.py}\quad（359行）\par}

\VerbatimInput{code/check_outer_timing_ablation.py}

\label{code:source:30:end}

\Needspace{7\baselineskip}

\subsection{最终执行与交付检查}

\label{code:source:31}

{\small\noindent\path{code/check_final_strategy_delivery.py}\quad（449行）\par}

\VerbatimInput{code/check_final_strategy_delivery.py}

\label{code:source:31:end}

\Needspace{7\baselineskip}

\subsection{原分问稿来源与结构检查}

\label{code:source:32}

{\small\noindent\path{code/check_manuscript_sources.py}\quad（61行）\par}

\VerbatimInput{code/check_manuscript_sources.py}

\label{code:source:32:end}

\Needspace{7\baselineskip}

\subsection{原版完整运行入口}

\label{code:source:33}

{\small\noindent\path{code/run_all.ps1}\quad（46行）\par}

\VerbatimInput{code/run_all.ps1}

\label{code:source:33:end}

\Needspace{7\baselineskip}

\subsection{母模型修订复算入口}

\label{code:source:34}

{\small\noindent\path{code/run_revision.ps1}\quad（153行）\par}

\VerbatimInput{code/run_revision.ps1}

\label{code:source:34:end}

\Needspace{7\baselineskip}

\subsection{原分问稿打包工具}

\label{code:source:35}

{\small\noindent\path{code/package_revision.py}\quad（65行）\par}

\VerbatimInput{code/package_revision.py}

\label{code:source:35:end}

\Needspace{7\baselineskip}

\subsection{原版诊断动态加载的冻结实现}

\label{code:source:36}

{\small\noindent\path{versions/v1_before_settlement_revision/code/optimize_dispatch.py}\quad（387行）\par}

\VerbatimInput{versions/v1_before_settlement_revision/code/optimize_dispatch.py}

\label{code:source:36:end}

\Needspace{7\baselineskip}

\subsection{XeLaTeX 编译入口}

\label{code:source:37}

{\small\noindent\path{compile.sh}\quad（36行）\par}

\VerbatimInput{compile.sh}

\label{code:source:37:end}

\Needspace{7\baselineskip}

\subsection{工作簿导出的执行目录副本}

\label{code:source:38}

{\small\noindent\path{outputs/final_strategy/runtime/export_final_workbooks.mjs}\quad（261行）\par}

\VerbatimInput{outputs/final_strategy/runtime/export_final_workbooks.mjs}

\label{code:source:38:end}

\Needspace{7\baselineskip}

\subsection{公式与版式专项审计}

\label{code:source:39}

{\small\noindent\path{outputs/logic_formula_review/verify_revision.py}\quad（88行）\par}

\VerbatimInput{outputs/logic_formula_review/verify_revision.py}

\label{code:source:39:end}

\Needspace{7\baselineskip}

\subsection{本代码附录生成器}

\label{code:source:40}

{\small\noindent\path{outputs/code_appendix/build_code_appendix.py}\quad（116行）\par}

\VerbatimInput{outputs/code_appendix/build_code_appendix.py}

\label{code:source:40:end}

\endgroup

